\chapter{Hordeolum (Stye) and Chalazion}
\index{Hordeolum (Stye) and Chalazion}
\label{sec:Hordeolum (Stye) and Chalazion}

\section{Overview}
Hordeolum (stye) and chalazion are common eyelid lesions affecting the meibomian or Zeis glands. A hordeolum is an acute bacterial infection, while a chalazion is a chronic sterile lipogranulomatous inflammation from gland obstruction. These conditions affect people of all ages and are among the most frequent ophthalmic complaints in primary care.
\section{Classification}

\begin{description}[font=\normalfont\bfseries,leftmargin=3em,labelwidth=2.5em]
  \item[Cause] Infectious / Inflammatory
  \item[Organ System] Eyes
  \item[Pathophysiology] Bacterial infection (hordeolum) or duct obstruction with lipogranulomatous inflammation (chalazion)
  \item[Duration] Acute (hordeolum) to Chronic (chalazion)
  \item[Onset] Sudden (hordeolum) or Gradual (chalazion)
  \item[Mode of Transmission] Autoinoculation
  \item[Clinical Course] Acute, self-limited or chronic, resolving
  \item[Severity] Mild
  \item[Extent of Disease] Localized
  \item[Age Group] All ages
  \item[Epidemiology] Very common; affects all ages; more frequent in adults
  \item[ICD-11 Code] 9A60.0 / 9A60.1
\end{description}

\section{Causes}
A hordeolum results from acute bacterial infection of the eyelid glands, most commonly Staphylococcus aureus, causing localized abscess formation. A chalazion develops when the meibomian gland duct becomes obstructed, leading to extrusion of gland contents into surrounding tissue and a chronic lipogranulomatous inflammatory response. Poor eyelid hygiene and chronic blepharitis are predisposing factors.

\section{Risk Factors}

\begin{itemize}[noitemsep]
  \item Chronic blepharitis
  \item Poor eyelid hygiene
  \item Contact lens use
  \item Recurrent styes
  \item Diabetes mellitus
  \item Seborrheic dermatitis
  \item Rosacea
  \item Immunosuppression
\end{itemize}

\section{Signs and Symptoms}

\subsection{Early symptoms}
\begin{itemize}[noitemsep]
  \item Early manifestations of hordeolum (stye) and chalazion may be subtle and nonspecific
\end{itemize}

\subsection{Common symptoms}
\begin{itemize}[noitemsep]
  \item \subsection{Common symptoms}
  \item \textbf{Common symptoms:}
  \item Hordeolum presents as an acute, painful, erythematous nodule at the eyelid margin with localized swelling and tenderness
  \item Chalazion presents as a firm, painless, slowly growing nodule within the eyelid, away from the margin
  \item External hordeolum points at the lid margin; internal hordeolum points on the conjunctival surface
  \item Tearing, foreign body sensation, and photophobia
  \item Blurred vision from mass effect on the cornea
  \item Pain or discomfort in the affected area
  \item Difficulty performing daily activities
\end{itemize}

\subsection{Less common symptoms}
\begin{itemize}[noitemsep]
  \item Atypical or less common presentations of hordeolum (stye) and chalazion may occur
\end{itemize}

\subsection{Emergency warning signs}
\begin{itemize}[noitemsep]
  \item Severe or worsening symptoms of hordeolum (stye) and chalazion require immediate medical evaluation
\end{itemize}
\section{Progression and Stages}
An untreated hordeolum typically suppurates and drains spontaneously within 3--5 days. A chalazion may persist for weeks to months, gradually enlarging before eventually resolving or requiring intervention. Chronic or recurrent lesions warrant evaluation for underlying systemic conditions.

\section{Complications}

\begin{itemize}[noitemsep]
  \item Preseptal or orbital cellulitis (rare, from extension of infection)
  \item Recurrent hordeola or chalazia
  \item Corneal astigmatism from large chalazion
  \item Scarring or lid deformity
  \item Reduced quality of life
  \item Emotional and psychological effects
\end{itemize}

\section{Diagnosis}

\begin{itemize}[noitemsep]
  \item Clinical diagnosis based on history and physical examination
  \item Slit-lamp examination to differentiate hordeolum from chalazion
  \item Lid eversion to assess internal lesions
  \item Biopsy for atypical or recurrent lesions to rule out sebaceous cell carcinoma
  \item Lab tests to check for underlying conditions
\end{itemize}

\section{Prevention}

\begin{itemize}[noitemsep]
  \item Maintain good eyelid hygiene with warm compresses
  \item Avoid touching or rubbing the eyes
  \item Treat underlying blepharitis or rosacea
  \item Replace eye makeup regularly
  \item Go for regular check-ups so any problems are caught early
\end{itemize}

\section{Treatment Options}

\begin{itemize}[noitemsep]
  \item Warm compresses applied 4--6 times daily for 10--15 minutes promote drainage
  \item Gentle lid massage after warm compresses
  \item Topical antibiotic ointments (bacitracin, erythromycin) for hordeolum
  \item Incision and curettage for persistent chalazion
  \item Intralesional corticosteroid injection for large chalazia
  \item Oral antibiotics for associated cellulitis
  \item Lifestyle changes and self-care
\end{itemize}

\section{Living with It}

\begin{itemize}[noitemsep]
  \item Follow your treatment plan as prescribed
  \item Keep up with regular appointments with your healthcare provider
  \item Adjust your daily habits as needed
  \item Reach out to family, friends, or support groups for help
  \item Keep learning about your condition and how to manage it
\end{itemize}

\section{Outlook}
Both hordeolum and chalazion carry an excellent prognosis with appropriate treatment. Most hordeola resolve within one week, while chalazia may take several weeks to months. Recurrence is common, especially in patients with underlying blepharitis or rosacea, but complications are rare with proper management.
